\subsection{直和的张量代数。自由模的张量代数。分次模的张量代数。}

设 A 是一个交换环，且 \(M = \bigoplus_{\lambda \in L} M_{\lambda}\) 是一族A-模的直和。由II，§3，第7号，命题7，通过对n归纳可知， \(T^{n}(M)\) 是子模的直和，这些子模是典范单射的像

\[
\mathbf {M} _ {\lambda_ {1}} \otimes \mathbf {M} _ {\lambda_ {2}} \otimes \dots \otimes \mathbf {M} _ {\lambda_ {n}} \rightarrow \mathsf {T} ^ {n} (\mathbf {M}) = \mathbf {M} ^ {\otimes n}
\]

相对于所有序列 \((\lambda_{i})\in\mathbf{L}^{n}\) 。将 \(M_{\lambda_{1}}\otimes M_{\lambda_{2}}\otimes\cdots\otimes M_{\lambda_{n}}\) 与这个像等同，可以看出 \(T(\mathbf{M})\) 是下列模的直和

\[
\mathbf {M} _ {\lambda_ {1}} \otimes \mathbf {M} _ {\lambda_ {2}} \otimes \dots \otimes \mathbf {M} _ {\lambda_ {n}}
\]

（其中 \(n\) 遍历 \(\mathbf{N}\) ，且对每个 \(n\) ， \((\lambda_t)\) 遍历 \(\mathbf{L}^n\) .

我们首先推导出以下结论：

定理1. 设A为交换环，M为自由A-模，且 \((e_{\lambda})_{\lambda \in L}\) 为 M 的一组基。则诸元素 \(e_{s} = e_{\lambda_{1}} \otimes e_{\lambda_{2}} \otimes \cdots \otimes e_{\lambda_{n}}\) ——其中 \(s = (\lambda_{1}, \ldots, \lambda_{n})\) 遍历 L 中元素的所有有限序列的集合，而 \(e_{\varnothing}\) 用来表示 T(M) 的单位元——构成 A-模 T(M) 的一组基。

该基的元素显然是齐次的，且乘法表由下式给出

\[
e _ {s} e _ {t} = e _ {s t}\tag{10}
\]

（其中 \(st\) 表示由 \(\mathbf{L}\) 中元素组成的、通过并置序列 \(s\) 和 \(t\) 而得到的序列（I，§7，第2号）。

可以看出，T(M) 的基 \((e_{s})\) 连同乘法法则 (10) 典范同构于集合 L 上的自由幺半群（I，§7，第2号），该同构通过把这个幺半群的每个字 s 映到元素 \(e_{s}\) 而得到。由此（§2，第7号）可知，自由模 M（其基的指标集为 L）的张量代数 T(M) 典范同构于 L 在 A 上的自由结合代数。特别地（§2，第7号，命题7），对于每个从 L 到 A-代数 E 的映射 \(f: L \to E\) ，存在唯一的 A-代数同态 \(\bar{f}: T(M) \to E\) ，使得 \(\bar{f}(e_{\lambda}) = f(\lambda)\) 。

注。上述结果同样可以由自由结合代数和张量代数的泛性质推出，利用II，§3，第7号，命题7的推论2。

推论. 若 M 是投射 A-模，则 T(M) 是投射 A-模。

M 是自由 A-模 N 的一个直因子（II，§2，第 2 号，命题 4），因此 T(M) 是 T(N) 的一个直因子（第 2 号）；由于 T(N) 是自由的（定理 1），这表明 T(M) 是投射的（II，§2，第 2 号）。

命题7. 设 \(\Delta\) 是一个交换幺半群， \(\mathbf{M}\) 是一个 \(\Delta\) 型分次 A-模，而 \((\mathbf{M}_{\alpha})_{\alpha \in \Delta}\) 是其分次。对于每个有序对 \((\alpha, n) \in \Delta \times \mathbf{N}\) ，设 \(\mathsf{T}^{\alpha,n}(\mathbf{M})\) 为诸子模 \(\mathbf{M}_{\alpha_1} \otimes \mathbf{M}_{\alpha_2} \otimes \dots \otimes \mathbf{M}_{\alpha_n}\) （属于 \(\mathsf{T}^n(\mathbf{M})\) ，满足 \(\sum_{i=1}^{n} \alpha_i = \alpha\) ）的（直）和；于是 \((T^{\alpha, n}(M))_{(\alpha, n) \in \Delta \times N}\) 是唯一的 \(\Delta \times N\) 型分次，它与 \(T(M)\) 上的代数结构相容、并在 \(M = T^1(M)\) 上诱导出给定的分次。

在本号开头已经看到， \(\mathsf{T}(\mathbf{M})\) 是诸 \(\mathsf{T}^{\alpha,n}(\mathbf{M})\) 的直和，而这是一个与代数结构相容的分次这一事实直接由定义得出。若 \((\mathsf{T}'^{\alpha,n})\) 是另一个 \(\Delta \times \mathbf{N}\) 型分次，它在 \(\mathsf{T}(\mathbf{M})\) 上与代数结构相容且满足 \(\mathsf{T}^{\alpha,1}(\mathbf{M}) = \mathsf{T}'^{\alpha,1}\) （对 \(\alpha \in \Delta\) ），则直接由定义可知，对所有 \(n \geqslant 1\) 以及所有 \(\alpha \in \Delta\) ，有 \(\mathsf{T}^{\alpha,n}(\mathbf{M}) \subset \mathsf{T}'^{\alpha,n}\) ；但由于 \(\mathsf{T}(\mathbf{M})\) 也是诸 \(\mathsf{T}^{\alpha,n}(\mathbf{M})\) 的直和，这意味着 \(\mathsf{T}'^{\alpha,n} = \mathsf{T}^{\alpha,n}(\mathbf{M})\) (II，§1，第8号，注记)。

